\section{Background and Related Work}
\label{sec:background}

\subsection{Time Series Forecasting Architectures}

Multivariate long-sequence forecasting has evolved through several architectural paradigms.
Early Transformer-based models such as Informer \cite{zhou2021informer} and Autoformer \cite{wu2021autoformer} apply attention over the \emph{temporal} axis, treating each time-step as a token.
This yields $\mathcal{O}(L^2)$ complexity and conflates temporal modelling with cross-variate modelling, since each token mixes all $C$ variates at a single time-step.

A decisive shift came with the observation that \emph{inverting} the token definition---making each variate a token---decouples variate correlation from temporal structure and reduces attention cost to $\mathcal{O}(C^2)$ when $C \ll L$.
iTransformer \cite{liu2024itransformer} operationalises this by embedding each variate's full $L$-step history into a single $D$-dimensional token, applying self-attention across the $C$ variate tokens for cross-variate mixing, and a per-variate FFN for feature transformation.
S-Mamba \cite{wang2024smamba} adopts the same variate-token layout but replaces the single backbone with a two-stage pipeline: a bidirectional Mamba block for variate correlation followed by an FFN block for temporal refinement.
DLinear \cite{zeng2023dlinear} challenges the necessity of complex operators entirely, showing that a simple decomposition plus linear projection is surprisingly competitive, motivating investigation into whether operator complexity should be matched to each task.
PatchTST \cite{nie2023patchtst} and TimesNet \cite{wu2023timesnet} explore patched temporal tokens and 2D convolutional representations respectively, further diversifying the design space.

Across these models a common pattern emerges: each design \emph{commits} to a specific operator for each architectural slot without empirical justification for that choice.
This motivates an operator-agnostic search over a unified operator family.

\subsection{Linear Input-Varying (LIV) Operators}

The LIV framework \cite{thomas2024star,massaroli2024liv} provides a unified algebraic view of sequence mixing operators.
Its central insight is to replace the \emph{fixed} weight matrix of a standard linear layer with a weight that is \emph{assembled from the input itself}.

In a standard linear layer the output is $Y = Wx$, where $W$ is a learned parameter independent of $x$.
In a LIV block, four learned linear projections of $x$ first produce feature vectors $B$, $C$, $S$, $V$:
\begin{equation}
  B = W_B x, \quad C = W_C x, \quad S = f_S(x), \quad V = W_V x,
  \label{eq:featurizer}
\end{equation}
where $W_B, W_C, W_V$ are plain \texttt{nn.Linear} layers (single matrix multiplies) and $f_S$ is either a linear layer with a learned bias or, for the input-dependent convolution case, a small two-layer MLP.
These are called the \emph{featurizer}.

The $L \times L$ token-mixing matrix $T$ is then \emph{assembled algebraically} from $B$, $C$, and $S$ — the four projections have distinct roles:

\begin{itemize}
  \item $\mathbf{b}_j$ and $\mathbf{c}_i$ (rows of $B$ and $C$) \textbf{dot-product to produce the scalar gate strength} of each entry $T_{ij}$.
        Their inner product $\langle\mathbf{c}_i, \mathbf{b}_j\rangle$ measures how strongly position $i$ should attend to position $j$.
  \item $S$ \textbf{provides the structural scalar} $M_{ij}$ that further modulates the gate strength: it encodes the \emph{pattern} (a decay factor for SSMs, a kernel value for convolutions, the constant 1 for attention, or absent for the diagonal case).
  \item $V$ \textbf{carries the payload}: it never enters the construction of $T$ at all.
        Once $T$ is assembled, the token-mixed output is obtained by $T \cdot V$, which routes the actual feature content.
\end{itemize}

Concretely, each entry of $T$ is:
\begin{equation}
  T_{ij} = \underbrace{\langle\mathbf{c}_i,\, \mathbf{s}_{ij} \odot \mathbf{b}_j\rangle}_{\text{gate strength modulated by structure}},
  \label{eq:tij}
\end{equation}
where $\mathbf{s}_{ij} \in \mathbb{R}^d$ is the \emph{structural modulator} derived from $S$ — for LOW\_RANK it is the all-ones vector; for SEMI\_SEPARABLE it is the element-wise cumulative product $\prod_{k=j}^{i-1}\mathbf{a}_k$ of transition vectors from $S$; for TOEPLITZ it is a scalar $K_{i-j}$ broadcast from $S$.

The full LIV forward pass is:
\begin{align}
  \mathrm{mixed}_i &= \sum_j T_{ij} \cdot \mathbf{v}_j, \quad \mathbf{v}_j = \text{$j$-th row of } V,
  \label{eq:token_mix} \\
  Y_i &= W \cdot \mathrm{mixed}_i,
  \label{eq:channel_mix}
\end{align}
where $W$ is a \emph{static} \texttt{nn.Parameter} — a fixed learned matrix independent of $x$.
The net computation is: (1) project $x$ to $B, C, S, V$ via linear layers; (2) assemble $T$ from $B, C, S$ algebraically; (3) mix content with $\mathrm{mixed} = T \cdot V$; (4) project with fixed $W$.

\subsubsection{Token Mixing Families}

Four choices of how $S$ constructs $\mathbf{s}_{ij}$ span the space of practical sequence operators.

\textbf{LOW\_RANK} — $\mathbf{s}_{ij} = \mathbf{1}/\sqrt{d}$ (constant, S unused):
\begin{equation}
  T_{ij} = \langle\mathbf{c}_i,\, \mathbf{b}_j\rangle / \sqrt{d}, \qquad T = CB^\top / \sqrt{d}.
\end{equation}
Every pair $(i,j)$ interacts with equal structural weight; row-wise softmax recovers multi-head attention \cite{vaswani2017attention}.
Cost $\mathcal{O}(L^2 d)$; full bidirectional context.

\textbf{SEMI\_SEPARABLE} — $\mathbf{s}_{ij} = \prod_{k=j}^{i-1}\mathbf{a}_k$ (S gives forget gate $\mathbf{a}_k$ per step):
\begin{equation}
  T_{ij} = \left\langle\mathbf{c}_i,\; \left(\prod_{k=j}^{i-1}\mathbf{a}_k\right) \odot \mathbf{b}_j\right\rangle \;\; (j \le i), \quad T_{ij} = 0 \;\; (j > i),
\end{equation}
where $\mathbf{a}_k \in (0,1)^d$ is the per-step forget gate (the $k$-th row of $S$), and $\mathbf{b}_j$ is the input contribution at position $j$.
The structural product $\prod_{k=j}^{i-1}\mathbf{a}_k$ decays the contribution of position $j$ exponentially as it propagates forward to position $i$; it is computed once via \texttt{torch.cumprod} and shared across all pairs without any per-element network call.
In this work we use CfC \cite{hasani2021cfc} as the sole SEMI\_SEPARABLE operator; its complementary gating specialises the featurizer as described in Section~\ref{sec:cfc_liv} below.

\textbf{TOEPLITZ} — $\mathbf{s}_{ij} = K_{i-j}\,\mathbf{1}$ (S provides scalar kernel, zero beyond window):
\begin{equation}
  T_{ij} = \langle\mathbf{c}_i,\, \mathbf{b}_j\rangle \cdot K_{i-j}, \quad K_{i-j} = 0 \;\text{if}\; |i-j| \ge k.
\end{equation}
$S$ is averaged across the feature dimension to yield a scalar kernel vector $K \in \mathbb{R}^k$ indexed by offset.
For a fixed kernel $K$ is an \texttt{nn.Parameter}; for the input-dependent variant, a two-layer MLP maps the global average of $x$ to $K$ — this is the only operator where $f_S(x)$ involves a nonlinear network \cite{poli2023hyena}.
Cost $\mathcal{O}(Lkd)$ for window size $k$.

\textbf{DIAGONAL} — $\mathbf{s}_{ij} = \mathbf{1}$ only when $i=j$, else zero (S unused):
\begin{equation}
  T_{ij} = \langle\mathbf{c}_i,\, \mathbf{b}_i\rangle\,\delta_{ij}, \qquad \mathrm{mixed}_i = \langle\mathbf{c}_i,\, \mathbf{b}_i\rangle\;\mathbf{v}_i.
\end{equation}
Only the diagonal of $T$ is non-zero; the scalar $\langle\mathbf{c}_i, \mathbf{b}_i\rangle$ gates how strongly position $i$'s value $\mathbf{v}_i$ passes through.
With $\mathbf{c}_i = \mathbf{1}$ and $\mathbf{b}_i = \mathrm{SiLU}(W_{\mathrm{gate}}\,x_i)$, this recovers the SwiGLU FFN \cite{shazeer2020glu} at $\mathcal{O}(Ld)$ cost with no cross-token interaction.

The four families trade off context range against cost:
\vspace{2pt}
\begin{center}
\small DIAGONAL $<$ TOEPLITZ $<$ SEMI\_SEPARABLE $\approx$ LOW\_RANK \quad (in context width)
\end{center}
\vspace{2pt}
while SEMI\_SEPARABLE and DIAGONAL share $\mathcal{O}(Ld)$ recurrent complexity versus $\mathcal{O}(L^2d)$ for LOW\_RANK.

\subsection{Recurrent Sequence Models}

Structured State-Space models (S4 \cite{gu2021s4}) and Mamba \cite{gu2023mamba} parameterise causal recurrences where a hidden state $h_t$ summarises all past inputs.
Both belong to the SEMI\_SEPARABLE family in the LIV taxonomy: $S$ supplies the per-step transition scalar $A_k$ and the T matrix unrolls the recurrence as a lower-triangular matrix of decayed contributions.
They are included here as the direct conceptual predecessors of CfC, which is the SEMI\_SEPARABLE operator used in this work.

\subsection{CfC-LNN and its LIV Representation}
\label{sec:cfc_liv}

CfC (Closed-form Continuous-time) networks \cite{hasani2021cfc}, derived from Liquid Time-Constant neurons \cite{hasani2022liquid}, solve continuous-time recurrent dynamics in closed form.
The standard CfC state update is:
\begin{equation}
  h_t = g_t \odot h_{t-1} + (1 - g_t) \odot \tilde{h}_t,
  \quad g_t = \sigma(W_g\,x_t + b_g),
  \label{eq:cfc_recurrence}
\end{equation}
where $g_t \in (0,1)^d$ is the \emph{forget gate} and $(1-g_t)$ is the \emph{input gate}, with the complementary constraint $g_t + (1-g_t) = 1$ enforcing a convex interpolation between the retained state and the new input $\tilde{h}_t = \tanh(W_{\mathrm{in}}\,x_t)$.

\subsubsection{CfC as a SEMI\_SEPARABLE LIV Block}

Unrolling Eq.~\eqref{eq:cfc_recurrence} from position $j$ to position $i$ shows that $h_i$ is a weighted sum of all past inputs:
\begin{equation}
  h_i = \sum_{j=0}^{i} \underbrace{\left(\prod_{k=j+1}^{i} g_k\right)}_{\text{survival from }j\text{ to }i} \odot \underbrace{(1-g_j) \odot \tilde{h}_j}_{\text{input absorbed at }j}.
  \label{eq:cfc_unrolled}
\end{equation}
This is exactly the SEMI\_SEPARABLE token mixing.
The LIV featurizer for CfC maps each input token to:
\begin{equation}
  \mathbf{b}_j = (1-g_j) \odot W_{\mathrm{in}}\,x_j, \quad
  \mathbf{a}_k = g_k = \sigma(W_g\,x_k), \quad
  \mathbf{c}_i = W_C\,x_i, \quad
  \mathbf{v}_j = W_V\,x_j,
  \label{eq:cfc_featurizer}
\end{equation}
so each entry of $T(x)$ becomes:
\begin{equation}
  T_{ij} = \left\langle W_C\,x_i,\;\; \left(\prod_{k=j}^{i-1} \sigma(W_g\,x_k)\right) \odot (1-\sigma(W_g\,x_j)) \odot W_{\mathrm{in}}\,x_j \right\rangle.
  \label{eq:cfc_tij}
\end{equation}
The scalar $T_{ij}$ thus encodes: \emph{how much of position $j$'s contribution}, modulated by its own input gate $(1-g_j)$, \emph{survives} through the chain of forget gates $g_{j}\!\cdot\ldots\cdot g_{i-1}$, \emph{as read out} by position $i$'s output projection $\mathbf{c}_i$.
The token-mixed output $\mathrm{mixed}_i = \sum_{j \le i} T_{ij}\,\mathbf{v}_j$ is a weighted sum of value vectors, where the weight of each past position decays with the cumulative product of all intermediate forget gates — recovering precisely the CfC hidden state $h_i$ in Eq.~\eqref{eq:cfc_unrolled}.

The complementary constraint $g_t + (1-g_t) = 1$ ensures that each gate pair $(A_k, B_k)$ partitions unity, stabilising gradient flow through the cumulative product and making CfC particularly well-suited to smooth continuous-valued signals such as exchange rates and temperature series.

\subsection{NSGA-II}

NSGA-II \cite{deb2002nsga2} is a fast elitist evolutionary algorithm for multi-objective optimisation.
It maintains a population ranked by Pareto dominance: an individual $a$ dominates $b$ if $a$ is no worse than $b$ on all objectives and strictly better on at least one.
Rank-0 individuals form the Pareto front.
Within each rank, crowding distance prioritises individuals in sparse objective-space regions, preserving diversity.
Parents are selected by tournament on (rank, crowding distance); offspring are produced by multi-point crossover and per-gene mutation.
Elitism preserves the top individuals across generations, guaranteeing monotone improvement of the Pareto front.

\subsection{TiDAR: Think-in-Diffusion Autoregressive Prediction}

Standard autoregressive (AR) forecasting generates the horizon sequentially at $\mathcal{O}(H)$ forward passes.
TiDAR \cite{tidar2024} proposes a hybrid training objective that enables parallel \emph{draft} prediction of all $H$ steps in a single pass, with optional AR refinement.

The $L+H$ tokens are processed jointly under a hybrid sparsity mask: lookback tokens attend causally among themselves; forecast positions are masked but attend bidirectionally to the full lookback.
The training loss combines a next-step AR objective on the lookback with a parallel prediction objective on the masked horizon:
\begin{equation}
  \mathcal{L}_{\mathrm{TiDAR}} = \frac{1}{1+\alpha}\bigl(\alpha \cdot \mathcal{L}_{\mathrm{AR}} + \mathcal{L}_{\mathrm{Diff}}\bigr).
  \label{eq:tidar_loss}
\end{equation}
At inference, three operating modes trade quality against speed: pure draft ($\mathcal{O}(1)$), partial-AR ($\mathcal{O}(k)$), and full-AR ($\mathcal{O}(H)$).

\subsection{Related Work}

\textbf{NAS for time series.}
Existing automated search for time series focuses on hyperparameter optimisation (learning rate, hidden size, patch length) rather than operator-family selection \cite{elsken2019neural}.
Cell-based NAS \cite{zoph2017neural} has been applied to RNN design but not to the modern SSM/attention/convolution pool.

\textbf{Unified sequence operators.}
Mamba \cite{gu2023mamba} reframes SSMs as selective state spaces; Hyena \cite{poli2023hyena} generalises convolution to long implicit kernels.
The LIV formalism \cite{thomas2024star,massaroli2024liv} subsumes these under a single token-mixing abstraction, making it a natural grammar for operator search.

\textbf{Parallel sequence decoding.}
Non-autoregressive Transformers \cite{gu2018nonautoregressive} and diffusion language models demonstrate that parallel generation is viable in discrete settings.
TiDAR \cite{tidar2024} extends this to a hybrid draft-then-refine framework suitable for continuous signals.
